\documentclass[11pt,letterpaper]{article}

% ── Packages ────────────────────────────────────────────────────────────
\usepackage[margin=1in]{geometry}
\usepackage{graphicx}
\usepackage{hyperref}
\usepackage{booktabs}
\usepackage{longtable}
\usepackage{enumitem}
\usepackage{xcolor}
\usepackage{fancyhdr}
\usepackage{titlesec}
\usepackage[backend=biber,style=numeric,sorting=none]{biblatex}
\addbibresource{references.bib}

% ── Page style ──────────────────────────────────────────────────────────
\pagestyle{fancy}
\fancyhf{}
\rhead{\small NOAA EMC — MCP/RAG Platform}
\lhead{\small RAG Knowledge Base Architecture}
\cfoot{\thepage}

% ── Title formatting ───────────────────────────────────────────────────
\titleformat{\section}{\large\bfseries}{\thesection}{1em}{}
\titleformat{\subsection}{\normalsize\bfseries}{\thesubsection}{1em}{}

\hypersetup{
    colorlinks=true,
    linkcolor=blue!60!black,
    citecolor=green!50!black,
    urlcolor=blue!70!black
}

% ── Document ────────────────────────────────────────────────────────────
\begin{document}

\begin{center}
{\LARGE\bfseries Retrieval-Augmented Generation (RAG)\\[4pt]
for the Global Workflow Knowledge Base}\\[12pt]
{\large A Researcher's Guide to How AI Finds Answers\\
in 200,000+ Documents}\\[20pt]
{\normalsize NOAA EMC — Environmental Modeling Center\\
MCP/RAG Platform Team\\[8pt]
May 2026}
\end{center}

\vspace{1em}
\hrule
\vspace{1em}

\tableofcontents
\newpage

% ════════════════════════════════════════════════════════════════════════
\section{What is RAG?}
% ════════════════════════════════════════════════════════════════════════

Retrieval-Augmented Generation (RAG) is a technique that gives an AI
assistant access to a curated library of documents \emph{at the moment
it answers your question}. Instead of relying solely on what the model
memorized during training, RAG retrieves the most relevant passages
from a knowledge base and feeds them to the language model as context.

Think of it this way: a language model without RAG is like a researcher
working from memory alone. A language model \emph{with} RAG is like
that same researcher who can instantly pull the right page from a
200,000-document library before answering.

\subsection{Why RAG Matters for Scientific Computing}

Large language models are trained on broad internet data. They know
general programming patterns and common scientific concepts, but they
do \emph{not} know:

\begin{itemize}[nosep]
    \item The specific structure of the NOAA Global Workflow
    \item Which Fortran subroutine handles vortex initialization in HAFS
    \item What environment variables \texttt{JGFS\_FORECAST} expects
    \item How ESMF couples the atmosphere to the ocean via NUOPC
    \item The correct \texttt{err\_chk} pattern for NCO production scripts
\end{itemize}

RAG bridges this gap by storing all of that domain knowledge in a
searchable format and retrieving it on demand.

\subsection{The Three Steps of RAG}

Every RAG system follows the same fundamental pattern:

\begin{enumerate}
    \item \textbf{Ingest} — Documents are broken into chunks, converted
    to numerical vectors (embeddings), and stored in a searchable
    database.
    \item \textbf{Retrieve} — When a question arrives, the system
    converts the question to the same vector format and finds the most
    similar document chunks.
    \item \textbf{Generate} — The retrieved chunks are passed to the
    language model as context, and the model generates an answer
    grounded in that evidence.
\end{enumerate}

The quality of a RAG system depends primarily on Step 1 (what you put
in) and Step 2 (how well you find it). Step 3 is handled by the
language model itself.

% ════════════════════════════════════════════════════════════════════════
\section{How RAG Connects to MCP}
% ════════════════════════════════════════════════════════════════════════

The Model Context Protocol (MCP) is the transport layer — it defines
how an AI assistant communicates with external tools and data sources.
RAG is one of those data sources.

In our system, the MCP server exposes tools like
\texttt{search\_documentation} and \texttt{explain\_with\_context}.
When the AI assistant needs domain knowledge, it calls these tools via
MCP. Behind the scenes, the tool performs the RAG retrieval (Step 2)
and returns the relevant passages. The assistant then uses those
passages to formulate its answer (Step 3).

\begin{center}
\fbox{\parbox{0.85\textwidth}{
\centering
\textbf{User Question} $\rightarrow$ \textbf{AI Assistant}
$\rightarrow$ \textbf{MCP Tool Call} $\rightarrow$
\textbf{RAG Retrieval} $\rightarrow$ \textbf{Answer with Evidence}
}}
\end{center}

The key insight: MCP handles the \emph{communication}, RAG handles
the \emph{knowledge}. They are complementary layers, not alternatives.

% ════════════════════════════════════════════════════════════════════════
\section{Embeddings: How Text Becomes Searchable}
% ════════════════════════════════════════════════════════════════════════

The core technology that makes RAG work is \emph{embedding} — the
process of converting text into a numerical vector (a list of numbers)
that captures its meaning.

\subsection{What is an Embedding?}

An embedding is a fixed-length list of numbers (a vector) that
represents the \emph{semantic meaning} of a piece of text. Two pieces
of text that mean similar things will have vectors that are close
together in this numerical space, even if they use completely different
words.

For example:
\begin{itemize}[nosep]
    \item ``FV3 dynamical core'' and ``cubed-sphere atmospheric solver''
    would have \emph{similar} vectors (same concept, different words)
    \item ``FV3 dynamical core'' and ``chocolate cake recipe'' would
    have \emph{distant} vectors (unrelated concepts)
\end{itemize}

This property — that meaning maps to proximity in vector space — is
what allows RAG to find relevant documents even when the user's
question doesn't use the exact same terminology as the source material.

\subsection{Amazon Titan Text Embeddings v2}

Our knowledge base uses \textbf{Amazon Titan Text Embeddings v2}, a
sentence transformer model provided by AWS through the Bedrock service.
Key characteristics:

\begin{itemize}[nosep]
    \item \textbf{1024 dimensions} — each text chunk is represented as
    a vector of 1024 numbers
    \item \textbf{8,192 token input limit} — can process chunks up to
    approximately 32,000 characters
    \item \textbf{Multilingual} — handles technical English, code
    snippets, and mixed content
    \item \textbf{Managed service} — runs on AWS infrastructure, no
    local GPU required
\end{itemize}

When we ingest a document, each chunk of text is sent to Titan via the
AWS Bedrock API. Titan returns a 1024-number vector that we store
alongside the original text in OpenSearch.

\subsection{Why We Switched to Titan}

Previously, the system used a locally-hosted sentence transformer model
(MPNet, 768 dimensions). The migration to Titan provided:

\begin{enumerate}[nosep]
    \item \textbf{Higher dimensionality} (1024 vs 768) — more capacity
    to distinguish subtle semantic differences
    \item \textbf{No local compute} — embedding generation runs on AWS,
    freeing the server for other work
    \item \textbf{Consistent quality} — AWS maintains and updates the
    model; we benefit from improvements without re-deploying
    \item \textbf{Scale} — Titan handles our 200,000+ document corpus
    without memory constraints
\end{enumerate}

The MPNet indices remain available as a fallback (the system maintains
both), but all new ingestion targets the Titan 1024-dimensional indices.

% ════════════════════════════════════════════════════════════════════════
\section{The Vector Store: OpenSearch}
% ════════════════════════════════════════════════════════════════════════

Embeddings need a home — a database optimized for finding the nearest
vectors to a query vector. We use \textbf{Amazon OpenSearch Service}
with its k-NN (k-nearest neighbors) plugin.

\subsection{How Vector Search Works}

When you ask a question:
\begin{enumerate}[nosep]
    \item Your question is embedded using the same Titan model
    \item OpenSearch finds the $k$ document chunks whose vectors are
    closest to your question's vector (using cosine similarity)
    \item Those chunks are returned as search results, ranked by
    similarity score
\end{enumerate}

This is fundamentally different from keyword search. A keyword search
for ``atmosphere model'' would miss documents that say ``atmospheric
solver'' or ``FV3 dynamical core.'' Vector search finds all of them
because their \emph{meanings} are close in embedding space.

\subsection{Our Index Structure}

The knowledge base is organized into five logical collections, each
stored as a physical OpenSearch index:

\begin{table}[h]
\centering
\begin{tabular}{lrl}
\toprule
\textbf{Index} & \textbf{Documents} & \textbf{Content} \\
\midrule
mdc-workflow-docs-titan1024 & 30,173 & Documentation from 53 URL sources \\
mdc-code-context-titan1024 & 90,135 & Fortran/Shell/Python code with context \\
mdc-jjobs-titan1024 & 751 & J-Job script headers and documentation \\
mdc-community-summaries-titan1024 & 2,113 & Graph-derived subsystem summaries \\
mdc-ee2-standards-titan1024 & 34 & EE2/NCO production coding standards \\
\bottomrule
\end{tabular}
\caption{OpenSearch indices in the Titan 1024-dim embedding space}
\end{table}

Total: over 123,000 documents in the Titan indices alone, with an
additional 85,000+ in the legacy MPNet indices for backward
compatibility.

% ════════════════════════════════════════════════════════════════════════
\section{The Unified Ingest Manifest}
% ════════════════════════════════════════════════════════════════════════

The \textbf{unified ingest manifest} is the single source of truth for
what documentation has been ingested into the knowledge base. It lives
at:

\begin{verbatim}
mcp_server_python/src/config/unified_manifest.json
\end{verbatim}

\subsection{What the Manifest Tracks}

For every documentation source in the knowledge base, the manifest
records:

\begin{itemize}[nosep]
    \item \textbf{name} — unique identifier (e.g., \texttt{mom6},
    \texttt{esmf-user-guide})
    \item \textbf{source\_type} — how the content is obtained
    (\texttt{url\_crawl}, \texttt{code\_parse}, \texttt{pdf\_download},
    etc.)
    \item \textbf{collection\_target} — which OpenSearch collection
    receives the documents
    \item \textbf{embedding\_profile} — which embedding model to use
    (\texttt{titan1024})
    \item \textbf{url} — the documentation website to crawl
    \item \textbf{doc\_count} — how many chunks are currently indexed
    \item \textbf{last\_ingested} — when the source was last
    successfully crawled
    \item \textbf{tier} — priority classification for ingestion ordering
    \item \textbf{enabled} — whether the source is active
\end{itemize}

\subsection{Source Types}

The manifest supports seven source types, reflecting the diverse nature
of the global workflow's documentation:

\begin{description}[style=nextline,leftmargin=2cm]
    \item[\texttt{url\_crawl}] Web documentation sites (ReadTheDocs,
    GitHub Pages, wikis). The crawler follows links within the site up
    to a configured page limit.
    \item[\texttt{pdf\_download}] PDF reference documents. Downloaded,
    text-extracted via pypdf, chunked, and embedded.
    \item[\texttt{on\_disk\_submodule}] Local RST/Markdown files from
    the global-workflow git submodule tree.
    \item[\texttt{code\_parse}] Source code (Fortran, Shell, Python)
    parsed at function boundaries with surrounding context.
    \item[\texttt{config\_parse}] Configuration files (Rocoto XML,
    bash key-value configs) parsed with structure-aware chunking.
    \item[\texttt{standards}] EE2/NCO coding standards from the
    nws-hpc-standards repository.
    \item[\texttt{community\_summary}] Graph-derived community
    summaries produced by the Leiden algorithm over the code graph.
\end{description}

\subsection{Tier System}

Sources are classified into five priority tiers that determine
ingestion order and operational importance:

\begin{table}[h]
\centering
\begin{tabular}{llp{7cm}}
\toprule
\textbf{Tier} & \textbf{Sources} & \textbf{Description} \\
\midrule
tier1\_critical & 9 & Core workflow, DA, coupling infrastructure
(ESMF, NUOPC, GSI, global-workflow, UFS-Utils) \\
tier2\_workflow & 5 & Workflow orchestration (ecFlow, Rocoto, wxflow,
pyflow, uwtools) \\
tier3\_models & 18 & Model components — atmosphere, ocean, ice, land,
chemistry, waves (MOM6, CICE, FV3, HAFS, MPAS, etc.) \\
tier4\_build & 15 & Build systems, I/O libraries (spack, NCEPLIBS,
wgrib2, Kokkos) \\
tier5\_standards & 6 & Coding standards and style guides (PEP8,
Fortran best practices, METplus) \\
\bottomrule
\end{tabular}
\caption{Documentation source tier classification}
\end{table}

\subsection{Gap Detection}

The manifest system includes automated gap detection that compares
declared sources against live OpenSearch data. For each collection, it
reports:

\begin{itemize}[nosep]
    \item \textbf{Coverage percentage} — actual documents vs. declared
    count
    \item \textbf{Never-ingested sources} — entries with no
    \texttt{last\_ingested} timestamp
    \item \textbf{Stale sources} — entries not refreshed in 30+ days
\end{itemize}

This allows operators to quickly identify which documentation sources
need attention without manually querying OpenSearch.

% ════════════════════════════════════════════════════════════════════════
\section{The Ingestion Pipeline}
% ════════════════════════════════════════════════════════════════════════

Getting documents from their source into the vector store involves a
multi-step pipeline.

\subsection{HTML Documentation Crawling}

For web-based documentation (the majority of sources), the pipeline:

\begin{enumerate}[nosep]
    \item Reads the manifest to identify enabled sources
    \item Crawls each URL, following internal links up to
    \texttt{max\_pages}
    \item Strips HTML to extract clean text content
    \item Splits text into semantic chunks (respecting paragraph and
    section boundaries)
    \item Sends each chunk to Amazon Titan for embedding (1024-dim
    vector)
    \item Indexes the chunk + vector + metadata into OpenSearch
    \item Updates the manifest with the new \texttt{doc\_count} and
    \texttt{last\_ingested} timestamp
\end{enumerate}

\subsection{PDF Document Ingestion}

For PDF-only reference documents (like the ESMF Fortran API reference),
a dedicated pipeline:

\begin{enumerate}[nosep]
    \item Downloads the PDF from the manifest URL
    \item Extracts text page-by-page using pypdf
    \item Chunks at 512-token windows with 64-token overlap between
    consecutive chunks
    \item Embeds each chunk with Titan (same 1024-dim space as HTML
    docs)
    \item Indexes with metadata including page number for traceability
\end{enumerate}

The 512-token chunk size with 64-token overlap ensures that:
\begin{itemize}[nosep]
    \item Each chunk is well within Titan's 8,192-token input limit
    \item Overlapping windows prevent information loss at chunk
    boundaries
    \item Retrieval returns coherent passages rather than fragments
\end{itemize}

\subsection{Code Ingestion}

Source code (90,000+ Fortran subroutines, shell scripts, and Python
modules) is ingested with function-boundary chunking — each chunk
corresponds to a complete function or subroutine with its surrounding
context (imports, module declarations, neighboring functions). This
ensures that code search results are self-contained and actionable.

% ════════════════════════════════════════════════════════════════════════
\section{Search: Putting It All Together}
% ════════════════════════════════════════════════════════════════════════

When a researcher asks a question through the MCP interface, the
\texttt{search\_documentation} tool performs a hybrid search:

\begin{enumerate}[nosep]
    \item \textbf{Semantic search (k-NN)} — the question is embedded
    with Titan and the nearest vectors are found in OpenSearch
    \item \textbf{Keyword search (BM25)} — traditional text matching
    catches exact terms the embedding might generalize over
    \item \textbf{Reciprocal Rank Fusion (RRF)} — results from both
    methods are merged using a rank-based scoring formula that
    preserves the best of each approach
    \item \textbf{Graph enrichment} (optional) — each result is
    annotated with the count of related entities in the code graph,
    helping identify which results connect to the broader system
\end{enumerate}

This hybrid approach means the system finds relevant results whether
the user asks in natural language (``how does the ocean couple to the
atmosphere?'') or with specific technical terms
(\texttt{ESMF\_FieldBundleCreate}).

% ════════════════════════════════════════════════════════════════════════
\section{Current State of the Knowledge Base}
% ════════════════════════════════════════════════════════════════════════

As of May 2026, the knowledge base contains:

\begin{itemize}[nosep]
    \item \textbf{209,277 total documents} across all indices
    \item \textbf{30,173 documentation chunks} in the Titan workflow-docs
    index (from 53 URL sources + 4 PDFs)
    \item \textbf{90,135 code chunks} covering Fortran, Shell, and
    Python from the global-workflow source tree
    \item \textbf{105,891 graph nodes} and \textbf{2.9 million
    relationships} in the Neptune code graph
    \item \textbf{64 declared sources} in the unified manifest (55
    url\_crawl + 9 non-URL)
\end{itemize}

The system is healthy (4/4 components operational) and serves queries
through both the AgentCore-hosted Python runtime and the local Node.js
gateway.

% ════════════════════════════════════════════════════════════════════════
\section{Summary}
% ════════════════════════════════════════════════════════════════════════

RAG transforms a general-purpose AI assistant into a domain expert by
giving it instant access to curated documentation at query time. The
key components are:

\begin{enumerate}[nosep]
    \item \textbf{Embeddings} (Titan 1024-dim) convert text to
    searchable vectors
    \item \textbf{OpenSearch} stores and retrieves those vectors
    efficiently
    \item \textbf{The manifest} tracks what's been ingested and
    identifies gaps
    \item \textbf{The ingest pipeline} keeps the knowledge base current
    \item \textbf{MCP} provides the communication layer between the AI
    and the RAG system
\end{enumerate}

The result: when a researcher asks about MPAS mesh generation or NUOPC
component coupling, the system retrieves the actual documentation —
not a hallucinated approximation — and grounds its answer in evidence.

\newpage

% ════════════════════════════════════════════════════════════════════════
\appendix
\section{Appendix: Complete Source Inventory}
\label{app:inventory}
% ════════════════════════════════════════════════════════════════════════

The following tables list every documentation source in the unified
ingest manifest as of May 20, 2026 (manifest version 9.0.0).

\subsection{Tier 1 — Critical (9 sources)}

\begin{longtable}{p{3.2cm}rp{7.5cm}}
\toprule
\textbf{Source} & \textbf{Docs} & \textbf{URL} \\
\midrule
\endhead
esmf-user-guide & 10,513 & earthsystemmodeling.org/\ldots/ESMF\_usrdoc/ \\
esmf-ref-pdf & 1,165 & earthsystemmodeling.org/\ldots/ESMF\_refdoc.pdf \\
esmc-ref-pdf & 125 & earthsystemmodeling.org/\ldots/ESMC\_crefdoc.pdf \\
nuopc-layer-ref & 248 & earthsystemmodeling.org/\ldots/NUOPC\_refdoc/ \\
nuopc-ref-pdf & 487 & earthsystemmodeling.org/\ldots/NUOPC\_refdoc.pdf \\
esmpy-pdf & 94 & earthsystemmodeling.org/\ldots/ESMPy.pdf \\
global-workflow & 236 & global-workflow.readthedocs.io/en/latest/ \\
ufs-utils & 129 & noaa-emcufs-utils.readthedocs.io/en/latest/ \\
gsi-user-guide & 87 & dtcenter.org/\ldots/gsi/docs/users-guide/ \\
\bottomrule
\end{longtable}

\subsection{Tier 2 — Workflow (5 sources)}

\begin{longtable}{p{3.2cm}rp{7.5cm}}
\toprule
\textbf{Source} & \textbf{Docs} & \textbf{URL} \\
\midrule
\endhead
ecflow & 338 & ecflow.readthedocs.io/en/latest/ \\
pyflow & 103 & pyflow-workflow-generator.readthedocs.io/ \\
rocoto & 74 & christopherwharrop.github.io/rocoto/ \\
uwtools & 152 & uwtools.readthedocs.io/en/latest/ \\
wxflow & 95 & wxflow.readthedocs.io/en/latest/ \\
\bottomrule
\end{longtable}

\subsection{Tier 3 — Models (18 sources)}

\begin{longtable}{p{3.5cm}rp{7cm}}
\toprule
\textbf{Source} & \textbf{Docs} & \textbf{URL} \\
\midrule
\endhead
catchem & 45 & ufs-community.github.io/CATChem \\
cdeps & 62 & escomp.github.io/CDEPS/ \\
cece & 38 & ufs-community.github.io/CECE \\
cice & 325 & cice-consortium-cice.readthedocs.io/ \\
cmaq & 403 & github.com/USEPA/CMAQ/wiki \\
cmeps & 0 & escomp.github.io/CMEPS/ \\
fms & 933 & github.com/NOAA-GFDL/FMS/wiki \\
fv3-docs & 851 & github.com/NOAA-GFDL/\ldots/wiki \\
gocart & 1,025 & geos-chem.readthedocs.io/ \\
hafs & 78 & hafsdoc.readthedocs.io/en/latest/ \\
jedi-docs & 107 & jedi-docs.readthedocs-hosted.com/ \\
land-da & 56 & land-da.readthedocs.io/en/stable/ \\
mom6 & 901 & mom6.readthedocs.io/en/main/ \\
mpas-atmosphere & 134 & www2.mmm.ucar.edu/projects/mpas/ \\
pyioda & 142 & jedi-docs\ldots/ioda/index.html \\
ufs-srweather-app & 189 & ufs-srweather-app.readthedocs.io/ \\
ufs-weather-model & 356 & ufs-weather-model.readthedocs.io/ \\
ww3-wiki & 2,282 & github.com/NOAA-EMC/WW3/wiki \\
\bottomrule
\end{longtable}

\subsection{Tier 4 — Build / Libraries (15 sources)}

\begin{longtable}{p{3.2cm}rp{7.5cm}}
\toprule
\textbf{Source} & \textbf{Docs} & \textbf{URL} \\
\midrule
\endhead
ccpp-techdoc & 214 & ccpp-techdoc.readthedocs.io/ \\
kokkos-api & 0 & kokkos.org/kokkos-core-wiki/ \\
kokkos-overview & 49 & kokkos.org/about/overview/ \\
nceplibs-bacio & 95 & noaa-emc.github.io/NCEPLIBS-bacio/ \\
nceplibs-bufr & 1,251 & noaa-emc.github.io/NCEPLIBS-bufr/ \\
nceplibs-g2 & 388 & noaa-emc.github.io/NCEPLIBS-g2/ \\
nceplibs-g2tmpl & 158 & noaa-emc.github.io/NCEPLIBS-g2tmpl/ \\
nceplibs-ip & 838 & noaa-emc.github.io/NCEPLIBS-ip/ \\
nceplibs-nemsio & 9 & noaa-emc.github.io/NCEPLIBS-nemsio/ \\
nceplibs-sfcio & 0 & noaa-emc.github.io/NCEPLIBS-sfcio/ \\
nceplibs-sigio & 1 & noaa-emc.github.io/NCEPLIBS-sigio/ \\
nceplibs-w3emc & 648 & noaa-emc.github.io/NCEPLIBS-w3emc/ \\
spack & 1,969 & spack.readthedocs.io/en/latest/ \\
spack-stack & 162 & spack-stack.readthedocs.io/ \\
wgrib2 & 366 & cpc.ncep.noaa.gov/\ldots/wgrib2/ \\
\bottomrule
\end{longtable}

\subsection{Tier 5 — Standards (6 sources)}

\begin{longtable}{p{3.8cm}rp{7cm}}
\toprule
\textbf{Source} & \textbf{Docs} & \textbf{URL} \\
\midrule
\endhead
fortran-best-practices & 876 & fortran-lang.org/learn/best\_practices/ \\
google-shell-style & 36 & google.github.io/styleguide/shellguide.html \\
metplus & 444 & metplus.readthedocs.io/en/latest/ \\
numpy-docstrings & 31 & numpydoc.readthedocs.io/ \\
pep8 & 42 & peps.python.org/pep-0008/ \\
upp & 198 & upp.readthedocs.io/en/latest/ \\
\bottomrule
\end{longtable}

\subsection{Non-URL Sources (9 sources)}

\begin{longtable}{p{3.8cm}p{2.5cm}rp{4.5cm}}
\toprule
\textbf{Source} & \textbf{Type} & \textbf{Docs} & \textbf{Description} \\
\midrule
\endhead
fortran-code-context & code\_parse & 77,613 & Fortran subroutines/functions \\
shell-code-context & code\_parse & 2,079 & Shell scripts (j-jobs, ush) \\
python-code-context & code\_parse & 8,922 & Python modules/utilities \\
rocoto-config & config\_parse & 187 & Rocoto XML workflow defs \\
expdir-configs & config\_parse & 1,297 & Experiment config files \\
global-workflow-rst & on\_disk & 1,759 & Local RST documentation \\
ee2-standards & standards & 34 & EE2/NCO coding standards \\
community-summaries & community & 2,113 & Graph-derived summaries \\
jjob-docs & jjob\_docs & 751 & J-Job script headers \\
\bottomrule
\end{longtable}

% ════════════════════════════════════════════════════════════════════════
\newpage
\printbibliography
% ════════════════════════════════════════════════════════════════════════

\end{document}
